\documentclass[11pt]{article}
% Shared print layout. The web's provisional design and MathJax are independent.
\usepackage[utf8]{inputenc}
\usepackage[T1]{fontenc}
\usepackage[spanish,es-nodecimaldot]{babel}
\usepackage{lmodern}
\usepackage{amsmath,amssymb}
\usepackage{graphicx}
\usepackage{geometry}
\geometry{a4paper,left=20mm,right=20mm,top=25mm,bottom=22mm,headheight=16pt,headsep=8mm}
\usepackage{fancyhdr,lastpage}
\usepackage{enumitem}
\usepackage{xcolor}
\usepackage{pgfplots}
\pgfplotsset{compat=1.18}
\usepgfplotslibrary{fillbetween}
\usepackage[hidelinks,pdfencoding=auto]{hyperref}
\hypersetup{pdflang=es-ES,pdfauthor={},pdfsubject={Matemáticas Aplicadas a las Ciencias Sociales II}}
\pagestyle{fancy}
\fancyhf{}
\fancyhead[L]{\small Matemáticas Aplicadas a las Ciencias Sociales II}
\fancyhead[R]{\small 2.\textsuperscript{o} Bachillerato}
\fancyfoot[L]{\footnotesize Inferencia - General | Edición revisada 1.0}
\fancyfoot[R]{\footnotesize Página \thepage\ de \pageref{LastPage}}
\renewcommand{\headrulewidth}{0.4pt}
\setlength{\parindent}{0pt}
\setlength{\parskip}{0.5em}
\setlist{itemsep=0.3em,topsep=0.3em}
\newcommand{\apartado}[1]{\subsection*{#1}}
\pgfplotsset{
  revision/.style={
    width=0.94\linewidth,height=7.5cm,
    axis lines=middle,grid=major,grid style={gray!18},
    tick label style={font=\small},label style={font=\small},
    xlabel={$x$},ylabel={$y$},samples=160,
    legend style={at={(0.5,-0.16)},anchor=north,draw=none,font=\small},
    clip=true
  }
}

\hypersetup{pdftitle={Inferencia - Examen general - Examen revisado}}
\begin{document}
\begin{minipage}{0.17\linewidth}\includegraphics[width=0.8\linewidth]{logo-al-andalus.png}\end{minipage}\hfill
\begin{minipage}{0.79\linewidth}{\Large\bfseries Inferencia: examen general}\\[0.4em]
Duración original: \textbf{55 minutos}.\end{minipage}

Apellidos y nombre: \hrulefill

Fecha: \hrulefill

Responde razonadamente y numera tus páginas. Se permite calculadora. Cada ejercicio incluye 0,5 puntos de expresión; no se añaden a su máximo.
Puntuación: \textbf{10 puntos}. Se conservan los cuatro ejercicios y sus apartados, sin añadir optatividad.
\section*{Ejercicio 1 \hfill 2,5 puntos}
 La Delegación de Tráfico de una ciudad desea estudiar la influencia del uso del teléfono móvil en los accidentes de tráfico. Elegida una muestra aleatoria simple de 250 accidentes registrados el año pasado, se observó que 90 de ellos se produjeron por distracciones debidas al uso del móvil.

\begin{enumerate}[label=(\alph*)]

  \item (1,5 puntos)  Determine un intervalo de confianza al 97\% para estimar la proporción de accidentes de tráfico debidos al uso del móvil mientras se conduce.

  

  \item (1 punto)  Usando la estimación anterior, calcule el tamaño mínimo que debe tener una muestra para estimar la proporción de accidentes con un error máximo del 5\% y un nivel de confianza del 99\%.

  

\end{enumerate}



\section*{Ejercicio 2 \hfill 2,5 puntos}
 Se desea estimar el porcentaje de jóvenes que utilizan una determinada red social. Para ello se escoge una muestra aleatoria simple de 500 jóvenes y de ellos 410 afirman utilizarla.

\begin{enumerate}[label=(\alph*)]
  \item (1,5 puntos)  Calcule el intervalo de confianza para la proporción de jóvenes que usa esa red social con un nivel de confianza del 95\%.

  

  \item (1 punto)  Manteniendo la proporción muestral, determine el tamaño mínimo de la muestra necesario para que, con un nivel de confianza del 97\%, el error máximo que se cometa al estimar la proporción de esa población sea inferior a 0.04.

  

\end{enumerate}



\clearpage
\section*{Ejercicio 3 \hfill 2,5 puntos}
 El peso de las ciruelas de una determinada variedad sigue una distribución Normal con media desconocida y desviación típica de 3 gramos. Se eligen al azar 25 ciruelas de esa variedad y se pesan, resultando un peso medio de 60 gramos.

\begin{enumerate}[label=(\alph*)]

  \item (1,5 puntos)  Calcule un intervalo al 95\% de confianza para estimar el peso medio de las ciruelas de esa variedad.

  

  \item (1 punto)  Calcule el tamaño mínimo de la muestra que se ha de tomar, para que al estimar el peso medio de esa variedad de ciruelas con un nivel de confianza del 99\%, el error cometido sea inferior a 1 gramo.

  

\end{enumerate}


\section*{Ejercicio 4 \hfill 2,5 puntos}
 El gasto que tienen los jóvenes durante un fin de semana es una variable aleatoria que sigue una distribución Normal de media \( \mu \) desconocida y desviación típica igual a 6 euros.

\begin{enumerate}[label=(\alph*)]

  \item (1,5 puntos)  Se toma una muestra aleatoria simple y se obtiene que el intervalo de confianza al 95\% para la media \( \mu \) es \( (24.47,\ 26.43) \). Calcule el valor de la media muestral y el tamaño de la muestra elegida.

  

  \item (1 punto)  Escogida otra muestra de tamaño 49 para estimar \( \mu \), calcule el error máximo cometido para esa estimación con un nivel de confianza del 97\%.

  

\end{enumerate}







\end{document}
